\documentclass[10pt,a4paper]{amsart}
\usepackage[margin=19mm]{geometry}
\usepackage{amsmath,amssymb,mathtools}
\usepackage[scr=rsfs,cal=euler]{mathalfa}
\usepackage{xfrac}
\usepackage[unicode,hidelinks,bookmarks=false]{hyperref}
\newcommand{\ie}{i.e.\ }
\newcommand{\cf}{cf.\ }
\newcommand{\resp}{respectively\ }
\newcommand{\Subsection}{\subsection}
\providecommand{\nobreakdash}{\nobreak\hbox{-}}
\setlength{\emergencystretch}{2em}
\multlinegap0pt
\usepackage{amssymb}
\usepackage{float}
\usepackage{bm}
\usepackage{tikz}
\newtheorem{theorem}{Theorem}
\newcommand{\Hh}{\mathbb{H}}
\renewcommand{\ge}{\geqslant}
\renewcommand{\le}{\leqslant}
\renewcommand{\tilde}{\widetilde}
\title{A reverse isoperimetric inequality in three-dimensional space forms}
\author{Kostiantyn Drach, Gil Solanes,, Kateryna Tatarko}
\date{Source: arXiv:2603.08132v1 (CC BY 4.0)}
\begin{document}
\maketitle

\noindent\emph{Source and licence.} A reverse isoperimetric inequality in three-dimensional space forms. Version arXiv:2603.08132v1 (CC BY 4.0), distributed by arXiv under the Creative Commons Attribution 4.0 International licence, http://creativecommons.org/licenses/by/4.0/, as recorded in the arXiv licence metadata for this exact version and shown on the abstract page https://arxiv.org/abs/2603.08132v1. Full licence notice: \texttt{licenses/arxiv-2603.08132-CC-BY-4.0.txt}.\par\smallskip

\noindent\textit{Editorial scope.} Counted unit: the article's theorem theorem (source0.tex lines 569--571). The article's complete original proof of it is delivered with it (source0.tex lines 573--654, one \texttt{proof} environment of 82 lines). No mathematics is written, repaired or re-derived: the statement and the proof are the article's own lines, in the article's own notation, and no retained line is substituted or re-flowed.  Retained uncounted as the source-local context the proof needs: lines 460--463; lines 477--483.  Nothing was omitted from any retained span, so the package carries no omission record.  The retained lines cite 2 works, whose entries are transcribed verbatim from the article's own bibliography source in the e-print and delivered with short editorial labels recorded in the item's reference mapping.  No retained line contains a figure environment, an image-inclusion command or drawing code, so no image is delivered and the package declares no asset dependency.  Editorial formatting only, in addition: an AMS \texttt{amsart} preamble that restates the article's own declarations and macros used by the retained lines, namely the environments \texttt{theorem} on separate counters, together with the standard packages those lines need.  The retained environments print in this document as Theorem 1 in order of appearance, whereas the article prints the same items with the numbering of its own full document.  The complete native source of the article as distributed in the arXiv e-print for this exact version is bundled unchanged as \texttt{sources/195984/source0.tex}.
\begin{equation}
    \label{Eq:Inrad}
    r(K) \ge r(L) \quad \text{ provided }\quad |\partial K| = |\partial L|,
\end{equation}

\begin{theorem}[Variation of the surface area for inner parallel polyhedra]
    \label{Thm:Variation}
    If $K \subset M^n(c)$ is a $\lambda$-convex polyhedron, then except for a finite number of $t_0 \in [0,r(K)]$,
    \[
    \left.\frac{d}{dt}\right|_{t=t_0} |\partial K_t| = -(n-1)\lambda(t_0)\cdot |\partial K_{t_0}| - 2\sum_{E \in \mathcal E(t_0)} \ell_E \cdot \tan\frac{\beta_E}{2}.
    \]
\end{theorem}

\begin{theorem}
    If $K \subset \Hh^2(-1)$ is a $\lambda$-convex body and $L \subset \Hh^2(-1)$ is a $\lambda$-convex lens such that $|\partial K| = |\partial L|$, then $|K| \ge |L|$. Moreover, equality is possible if and only if $K$ is a $\lambda$-convex lens.
\end{theorem}

\begin{proof}
    As in the proof of the Main Theorem, we can focus only on the case when $K$ is a $\lambda$-convex polygon. Let $e_0, \ldots, e_{m-1}$ ($m \ge 2$) be sides of $K$ and $\beta_i$ be the angle between the outer normals at the common vertex of $e_i$ and $e_{i+1}$ (indices are taken modulo $m$).    
    Likewise, let $\beta^*$ be the angle between the outer normals to the sides at the vertices of $L$. 

   We have that 
    \begin{equation}
        \label{Eq:Angles}
        \beta_i \le \beta^*, \quad \forall i.
    \end{equation}
    Indeed, for a fixed $i$, by Blaschke's rolling theorem, $K$ is contained in the set $\tilde{L}$ (compact or not), which is formed as the intersection of two supporting $\lambda$-balls that form sides $e_i$ and $e_{i+1}$. Thus, $|\partial K| \le |\partial \tilde{L}|$ by monotonicity and $|\partial L| \le |\partial \tilde{L}|$. We note that $\beta_i$ is the angle between the outer normals to the sides at a vertex of~$\tilde{L}$. 
    Since the perimeter of $\tilde L$ is non-decreasing with respect to $\beta_i$, we get that $\beta_i \le \beta^*.$ 

    
    By the two-dimensional Gauss--Bonnet theorem applied to $K$, we conclude
    \begin{equation}\label{Eq:GB2}
    2\pi = -|K| + \lambda |\partial K| + \sum_{i=0}^{m-1} \beta_i.
    \end{equation}
    Similarly, for the lens $L$, the Gauss--Bonnet formula gives
    \[
    2\pi = -|L| + \lambda |\partial L| + 2 \beta^*.
    \]
    Therefore, we conclude that 
    \begin{equation}
        \label{Eq:Cond}
        \sum_{i=0}^{m-1} \beta_i \le 2\beta^*, \quad \quad \text{whenever }\quad |\partial K| = |\partial L| \quad \text{and} \quad |K| \le |L|. 
    \end{equation}    

As in the proof of the Main Theorem, let $K_t$ be the inner parallel body for $K$ at distance $t \in [0, r(K)]$, and similarly $L_t$. Note that $K_t$ is a $m_t$-sided polygon, with $m_t \le m$. Recall that $r(L) \le r(K)$ provided that $|\partial K| = |\partial L|$ (see \eqref{Eq:Inrad}).

Since for every $t$, both $K_t$ and $L_t$ are $\lambda_t$-convex (for some $\lambda_t$ that can be computed explicitly depending on $t$ and the curvature of the ambient space), we can specify \eqref{Eq:Cond} for every $t_0 \in [0,r(L)]$:
\begin{equation}
        \label{Eq:Cond2}
        \sum_{i=0}^{m_{t_0}-1} \beta_{i,t_0} \le 2\beta^*_{t_0}, \quad \quad \text{whenever }\quad |\partial K_{t_0}| = |\partial L_{t_0}| \quad \text{and} \quad |K_{t_0}| \le |L_{t_0}|. 
    \end{equation}    
Here, $\beta_{i,t}$, $i \in \{0, \ldots, m_t - 1\}$, and $\beta_{*,t}$ are the corresponding angles of $K_t$ and $L_t$.

By Theorem~\ref{Thm:Variation},
\begin{equation}
\label{Eq:Der}
\begin{aligned}
    \left.\frac{d}{dt}\right|_{t=t_0}|\partial K_t| - \left.\frac{d}{dt}\right|_{t=t_0}|\partial L_t| &= 2 \tan \frac{\beta^*_{t_0}}{2} - 2 \sum_{i=0}^{m_{t_0}-1} \tan \frac{\beta_{i,t_0}}{2}
\end{aligned}
\end{equation}
for every $t_0 \in [0, r(L)]$ such that $|\partial K_{t_0}|= |\partial L_{t_0}|$.



Observe that $2x^{-1}\tan(\frac{x}2)$ is a monotone increasing function for $x \in [0, \pi)$. Assuming $|\partial K_{t_0}| = |\partial L_{t_0}|$ and $|K_{t_0}| \le |L_{t_0}|$, we can combine \eqref{Eq:Angles} and \eqref{Eq:Cond2} to conclude that 
\[
\sum_{i=0}^{m_{t_0}-1} \tan \frac{\beta_{i,t_0}}{2} =\sum_{i=0}^{m_{t_0}-1} \frac{\beta_{i,t_0}}{2} \cdot \frac{2}{\beta_{i,t_0}} \tan \frac{\beta_{i,t_0}}{2} \le \frac{2}{\beta^*_{t_0}} \tan \frac{\beta^*_{t_0}}{2} \sum_{i=0}^{m_{t_0}-1} \frac{\beta_{i,t_0}}{2} \le 2\tan \frac{\beta^*_{t_0}}{2}.
\]
Together with \eqref{Eq:Der}, this gives
\begin{equation}
    \label{Eq:Der2}
    \begin{aligned}    
    \left.\frac{d}{dt}\right|_{t=t_0}|\partial K_t| \ge  \left.\frac{d}{dt}\right|_{t=t_0}|\partial L_t| \quad &\text{for every }\quad t_0 \in [0, r(L)] \quad \\
    &\text{ such that } \quad |\partial K_{t_0}|= |\partial L_{t_0}| \quad \text{and} \quad |K_{t_0}| \le |L_{t_0}|. 
    \end{aligned}
\end{equation}
Moreover, the first inequality is strict if the last inequality $|K_{t_0}| \le |L_{t_0}|$ is also strict. 

Towards the contradiction, assume that $|K| <|L|$. Then we can apply \eqref{Eq:Der2} to conclude that 
\[
\left.\frac{d}{dt}\right|_{t=0}|\partial K_t| >  \left.\frac{d}{dt}\right|_{t=0}|\partial L_t|.
\]
Hence, the graph of the function $t \mapsto |\partial L_t|$ lies below the graph of the function $t \mapsto |\partial K_t|$ in some right neighborhood of $0$. If this is true for all $t \in [0,r(L)]$, we conclude that $|K| \ge |L|$, a contradiction. Therefore, there must exist the smallest $t_1 > 0$ such that 
\begin{equation}
\label{Eq:Der3}    
|\partial K_{t_1}| = |\partial L_{t_1}| \quad \text{ and }\quad \left.\frac{d}{dt}\right|_{t=t_1}|\partial K_t| \le   \left.\frac{d}{dt}\right|_{t=t_1}|\partial L_t|
\end{equation}
and the graph of $t \mapsto |\partial L_t|$ goes above the graph of $t \mapsto |\partial K_t|$ in a right neighborhood of $t_1$. By the choice of $t_1$, we know that 
\[
\int_0^{t_1}|\partial K_t|dt \ge \int_0^{t_1}|\partial L_t|dt. 
\]
Thus, using the assumption $|K| < |L|$,
\[
|K_{t_1}| = |K| - \int_0^{t_1}|\partial K_t|dt < |L| - \int_0^{t_1}|\partial L_t|dt = |L_{t_1}|.
\]
Therefore, at $t = t_1$, we have $|K_{t_1}| < |L_{t_1}|$ and $|\partial K_{t_1}| = |\partial L_{t_1}|$. Hence, by \eqref{Eq:Der2}, we must have $\left.\frac{d}{dt}\right|_{t=t_1}|\partial K_t| >  \left.\frac{d}{dt}\right|_{t=t_1}|\partial L_t|$. This is a contradiction to the inequality in \eqref{Eq:Der3}. This is a final contradiction that justifies $|K| \ge |L|$.

The equality case follows as in \cite{DT} (using the rigidity of the inequality in \cite{Dr}).
\end{proof}

\begin{thebibliography}{99}
\bibitem[DT]{DT} K. Drach, K. Tatarko, \emph{Reverse isoperimetric problems under curvature constraints}, (2023), \url{https://arxiv.org/abs/2303.02294}.
\bibitem[Dr]{Dr} K. Drach, \emph{Inradius estimates for convex domains in 2-dimensional Alexandrov spaces}, Anal. Geom. Metr. Spaces \textbf{6} (2018), no. 1, 165--173.
\end{thebibliography}
\end{document}
